\section{Evidence and completion status}
This appendix separates protocol details, equal-fact comparisons, local
resource accounting, exploratory MAG observations and the complete frozen
Amazon comparison families. Each has a different evidential role: a
reconstructed action is not a new labeled evaluation, a local timing is
not a deployed cost, and a frozen comparison family does not account
for earlier choices in the research program.

\section{Claim-to-evidence trace and reproducibility boundary}\label{app:trace}
The local analysis tree contains archived protocols, request and action
snapshots, source hashes and validation reports. A reader can verify
arithmetic without rerunning models, but cannot independently reproduce
end-to-end behavior from the PDF alone. Table~\ref{tab:trace} specifies
what evidence should be inspected for each claim. Paths are workspace
relative; access to a path is not a license to redistribute its contents.
This manuscript's English wording must remain consistent with the
more detailed Chinese working draft and the read-only evidence validator.

\begin{table}[h]
\centering
\caption{Audit targets under \texttt{output/analysis/}. Numeric IDs identify
the corresponding numbered report; H, I and queue IDs denote frozen
experiment directories. A local validation is not external replication.}
\label{tab:trace}\small
\begin{tabular}{p{5cm}p{2.7cm}p{4.5cm}}
\toprule
Claim & Evidence ID & What it checks \\
\midrule
H 18D comparison & H 20260922 & Protocol, actions, 384 hits, ten tests \\
I 20D comparison & I 20260923 & Different method, 768 hits, 15 tests \\
Same-fact rule replay & Report 66 & H/I label-blind action identity \\
Six graph columns & Report 67 & Scalar tolerance and input hashes \\
H text learner & Report 68 & Post hoc fit, gate, actions \\
Local resource scopes & Report 70 & Separate timed boundaries \\
Unscored ID queue & Queue v0 & Zero-access original protocol \\
Stopped launch & Launch v3 & No query/call clearance \\
Human review design & Report 65 & Prepared rubric, zero labels \\
\bottomrule
\end{tabular}
\end{table}

H and I are the directories \texttt{multiview\_transfer384\_20260922/}
and \texttt{disjoint\_transfer768\_20260923/}. Numbered IDs refer to
the matching Markdown report prefix in the analysis directory. Queue v0
and Launch v3 refer to \texttt{amazon\_candidate\_cohort\_20260925\_v0/}
and \texttt{amazon\_prospective\_launch\_20260925\_v3/}, respectively.

The read-only command \texttt{perl -w paper\_evidence\_check.pl}, from the
repository root, checks H/I summary and paired rows against frozen result
files, original source and input hashes for the H text pilot, candidate-ID
disjointness and zero-access status, and the two distinct H resource
benchmarks. It does not fetch the official dataset, rebuild the BM25/E5
indices, reproduce service outputs, retrain Qwen, run a blinded reviewer,
or prove that outside access to the candidate suffix never occurred.
Some raw records and model files are locally archived under an ignored
analysis tree; the project has not delivered a cleared public package.
The generated PDF and all third-party styles are reading artifacts, not
a substitute for inspectable authorized sources and exact version pins.

\paragraph{What a third-party reproduction would have to specify.}
At minimum: the STaRK revision and splits, full-text and graph schema,
hashes of every source/input/model, the lexical query tokenizer,
BM25/E5 initial retrieval and tie ordering, exact prompt and model IDs,
per-request caps and error semantics, entity-linking options/weights,
Qwen scoring revision, fit IDs and hyperparameters, 20 candidate IDs for
every query, action serialization, and a rule for missing responses.
A new run must specify whether API usage is billed and which units are
returned by the provider. A local syntax or checksum check does not
validate the truth of the benchmark labels or the legitimacy of
redistributing the underlying responses. Independent execution may
produce different API responses if the pinned service implementation
is unavailable; in that case the difference must be reported rather
than silently substituted with another model.

\section{Human semantic review is a separate estimand}\label{app:human}
The official synthetic qrels measure agreement with the benchmark
answer construction, not independent assessment of product role,
compatibility or intended purchase. The already prepared H/I blind
review draws per cohort 48 graph-versus-top-link \emph{disagreement}
queries and 24 queries on which they agree: 144 tasks in all. H has
73/384 differing actions and I has 197/768; the sample is chosen from
these exposed TRAIN cohorts, so any eventual semantic result remains
post hoc and cannot transform them into an untouched confirmation.
Each task displays at most the two products selected by the studied
policies. It does not expose all 20 retrieved products and cannot
estimate full-pool human recall.

Two real independent reviewers should receive separately shuffled
blinded forms containing the same source query, chosen-product text,
tags and recorded reference-edge facts. They should not be shown
policy names, frozen benchmark labels, previously computed outcomes
or the private mapping from displayed options to policy identities.
Each must label product role match, explicit textual conflict,
relational-condition support and overall suitability for \emph{each}
displayed option, using yes/no/uncertain with a cited justification.
An also-buy/also-view edge is evidence of co-occurrence, not proof
of semantic compatibility; missing graph evidence is not a verified
negative. Separating the two reviewers' access from the workspace
containing the private mapping is an operational requirement, not
something that a validator script can ensure.

Before unblinding, validate completeness and the two identities, archive
first-pass judgments, report uncertain rates and categorical agreement,
and perform a separately recorded adjudication. Only then join the
graph-versus-rule actions to suitability judgments, reporting fixes,
harms, both-yes, both-no and unresolved separately for H and I. Treat
unresolved paired differences conservatively with their possible
$[-1,1]$ contribution rather than dropping them. The agree-action
tasks serve as rubric and baseline checks, not evidence of a policy
difference. The disagreement-stratified sample is not IID from
all benchmark queries. These are analysis instructions; as of this
draft \emph{no human annotations have been completed}. Even completed
judgments cannot replace a new, independently frozen algorithmic
evaluation.

\section{Prospective study: stop conditions, not outcomes}\label{app:future}
A later preflight froze 384 official TRAIN \emph{IDs} selected from the
sorted-ID permutation at positions 4672--5055. Its immutable v0 record
sets paid calls to zero and forbids opening the new query text and
answers. It proposes a graph policy, the separately H-designed
structured-text learner and top-link on identical 20-candidate pools,
with a label-blind minimum net action-disagreement threshold of five
against \emph{both} controls. This gate is necessary for a meaningful
Hit@1 contrast but cannot predict the sign of the contrast. Passing
it would not retrospectively remove H's influence on text-arm design.

User authorization later specified a task-bounded maximum of 1,152
external requests, conditional on a complete protocol and on the
action gate before reading answers. The archived executor caps the
attempt count and per-request output tokens and verifies observed
service IDs, but records input usage and a provider cost field only
\emph{after} each call. It does not enforce a cumulative monetary
hard stop before dispatch, and the cost field's billing unit and an
explicit finite monetary ceiling remain unresolved. The separate
launch preflight therefore returned STOP: no new-query text, answers
or paid models were accessed by that attempt. Neither a user's
allowance estimate nor a locally frozen ID list is proof that the
cost-control or provenance conditions have been met.

To turn that queue into a real prospective experiment, one must
first resolve remaining possible exposure outside known local
metadata, verify the original service models without disclosing
credentials, bind the action producer and its input/model files to
an immutable protocol, implement enforceable token and monetary
limits, and specify abstention and missing-call treatment. Only
then may the approved query-only action generation run. The
independently fitted equal-fact text arm should use the same cached
edge/link facts, candidate pool and supervision and declare whether
the target is pure quality at fixed facts or quality per measured
resource. If full cold/warm lookup, comparable index construction
and updates, token counts, latency, RAM/VRAM and amortized costs
cannot be measured, that run cannot establish deployment superiority.
Scoring any answers remains conditional on an archived label-blind
action file that passes both predeclared gates. An unfavorable outcome
must be reported without extending the queue or changing the policy.

This leaves four independent completion criteria for a mature paper:
(i) an eligible prospectively scored comparison under correct
model/budget guards, (ii) genuinely matched end-to-end resource
measurements or an explicitly quality-only claim, (iii) actual
independent human semantic judgments with reviewer separation, and
(iv) author-approved reproduction rights and a redistributable
evidence bundle. None of these is manufactured by typesetting.

 \section{Implementation, data and frozen decisions}\label{app:implementation}
 \paragraph{Dataset and candidate universe.}
 The local STaRK-Amazon snapshot contains 957,192 products, 1,035,542
 nodes, and 5,382,103 stored directed edges. Stored edge counts depend on
 the snapshot's counting convention; they need not match a graph statistic
 reported for another revision. The full graph contains also-buy,
 also-view, brand, category and color edges, but the support proxy in this
 paper uses the first two types. The official synthetic query split contains
 5,910 TRAIN, 1,548 validation and 1,642 TEST queries, plus 81 human queries.
 Synthetic query generation uses both text and graph information; that
 design may favor relational features relative to naturally sampled tasks.
 Product texts, graph structure, queries and model pretraining corpora can
 overlap even when evaluation \emph{query IDs} are disjoint. A product not
 in the labeled answer set is a benchmark non-hit, not necessarily a human
 verified semantic negative.

 \paragraph{Retrieval, ranking and failure handling.}
 Frozen full-corpus BM25 and E5 ranks are joined using RRF with rank
 constant 60, and the first 20 candidates are retained. This candidate
 retrieval rank is distinct from the two subsequent LLM ranks of those
 20 candidates. Both the learned fusion and the top-link rule use the same
 frozen candidate order as the ultimate tie-breaker. The two service IDs
 recorded for H are \texttt{gpt-5.6-sol} and \texttt{gpt-6-astra}; the
 identity check uses reported service usage rather than model self-report.
 The reference plan is a separate shared request to the latter service,
 limited to three literal product mentions in the query. An exact title
 match is attempted before dense E5 link proposals. Each LLM sees at most
 512 \texttt{o200k\_base} tokens of product text plus capped tags, but
 the E5 retrieval truncation and Qwen effective contexts are not all
 identical token windows. No claim of equal input length across encoders
 follows from sharing product identities.

 A valid ranking contains every candidate exactly once. No ranking is
 silently repaired by appending absent products, and no parsed failure is
 removed from the query denominator. All arms in a cohort use the same
 usable service outputs. H had valid reference plans for all 384 queries
 and 765 valid ranking responses out of 768; I had valid plans for all 768
 and 1,533 valid rankings out of 1,536. The original 18D fusion on H
 scores 7,680 product/query pairs locally with a frozen Qwen3-Reranker-0.6B
 yes-versus-no logit difference. Calls to obtain ranks and reference plans
 are shared by top-link, but the rule does not require Qwen prediction or
 learned-model inference to make its decision. A same-query prediction
 cost measured after all upstream features are cached cannot be used as
 its deployment cost.

 \paragraph{Six rank and six nonrelational columns.}
 For each valid model $m$, write $r_m(v)$ for its rank and $M=|\mathcal M_q|$.
 The rank block actually used on the fixed $K=20$ pool is
 \begin{align}
 \phi_R(q,v)=\big[& (\min_m r_m(v)-1)/19,
  (\max_m r_m(v)-1)/19,\\
 & (M^{-1}\!\sum_m r_m(v)-1)/19,
  (\max_m r_m(v)-\min_m r_m(v))/19,\nonumber\\
 & M^{-1}\!\sum_m 1/r_m(v),
  M^{-1}\!\sum_m\mathbf{1}\{r_m(v)=1\}\big].\nonumber
 \end{align}
 The nonrelational block holds normalized initial retrieval rank;
 query/title and query/text term overlaps; a capped normalized product-text
 length; the frozen Qwen relevance logit difference $z(q,v)$; and
 $z(q,v)-\max_{u\in C_q}z(q,u)$. The Qwen instruction distinguishes a
 requested purchase from a reference item already owned by a user, but
 its score is not a verified role or calibrated probability. The text
 ablation is $[\phi_R;\phi_T]$, the graph method appends $\phi_G$,
 and shuffle preserves the six graph columns' marginal values while
 moving the candidate association as one block within each query.
 Shuffling is controlled by a fixed seed and query identifier. Only one
 prespecified shuffle realization was evaluated; even a successful
 paired comparison does not estimate variance across many shuffles.

 \paragraph{Six graph columns and policy.}
 For options $R_{qj}$, the similarity softmax is
 \begin{equation}
  w_{qjt}=\frac{\exp(s_{qjt}/0.05)}
  {\sum_{u\in R_{qj}}\exp(s_{qju}/0.05)},\qquad
  p_q(h)=\prod_jw_{qj,h_j}.
 \end{equation}
 Exact-title option ties instead receive uniform weights. If $A(v,t)$
 denotes the observed existence of an also-buy or also-view edge in
 either direction, $U_q(v,h)=J_q^{-1}\sum_j A(v,h_j)$ for nonzero $J_q$.
 We then take $\phi_G=[\mathbb E_{\rm unif} U,\min U,\max U,
 \operatorname{Std}_{\rm unif}U,\mathbb E_{p_q}U,
 \operatorname{Std}_{p_q}U]$. These moments flatten typed directed
 records into a coarse undirected existence indicator for scoring.
 The fact serialization used by the equal-fact controls retains the
 underlying edge types, directions and duplicates instead. This distinction
 matters: the six coarse values are not a semantics oracle. If no mention
 is identified, the whole graph block is zero.

 For top-link let $h_j^*=\arg\max_{t\in R_{qj}}w_{qjt}$, resolving a link
 tie by frozen option order. Its lexicographic action is
 \begin{equation}
 \pi_{\rm top}(q)=\arg\max_{v\in C_q}^{\rm lex}
  \left(U_q(v,h^*),\sum_{m\in\mathcal M_q}\frac{1}{60+r_m(v)},
  -\operatorname{rank}_{C_q}(v)\right).
 \end{equation}
 The second and third entries matter only after support ties. This is
 not a ``weak text-only'' baseline: it shares links and graph data with
 the graph fusion, but uses a different decision constraint. Max-world
 support replaces the first entry by $\max_hU_q(v,h)$, with the same
 two-rank tie-breaker. Its 276 H and 570 I hits are below top-link but
 above either learned graph method on I. A model trained without the
 top-link support column cannot be declared to have learned the
 support-first policy merely because the underlying facts were available.

 \subsection{Equal-fact controls: what is and is not tested}\label{app:equalfacts}
 \paragraph{Rule and feature replays.}
 The post hoc rule control serializes each query's frozen candidate order,
 rankings, reference option IDs, option weights, and all recorded connected
 candidate/option edge IDs, types, directions and duplicates as UTF-8
 structured text. An independent parser then selects each top link,
 recomputes its support and applies the same lexicographic RRF tie-break.
 It agrees on 384/384 H actions and 768/768 I actions; only 68 H and 138 I
 queries contain a named reference. The serialized per-query size spans
 215--8,480 H bytes and 228--3,891 I bytes. These counts specify \emph{bytes},
 not a token budget for a model. A second label-blind replay enumerates
 all linked-option worlds and reconstructs the original graph block
 to absolute error at most $10^{-10}$ for every one of the 46,080 H and
 92,160 I scalars (queries times 20 candidates times six columns).
 That construction can be performed in either a graph-oriented or
 text-parsed program. It does not establish that an off-the-shelf
 language model, given the same tokens, will learn an equivalent policy.

 \paragraph{An H-informed learned text arm.}
 To avoid calling a mere reserialization of six graph columns a learned
 text baseline, a separate H development experiment feeds the losslessly
 serialized \emph{facts} to a parser that constructs 18 lossy lexical
 statistics per candidate. These are combined with the original 12
 nonrelational features and fitted with the same 768 fitting queries,
 query groups, labels, LambdaRank settings and tree count. Original text
 and original graph models were retrained in the same CPU session and
 replayed their frozen H actions. The new text model makes 278/384
 H hits, against 278/384 graph and 279/384 top-link. Relative to graph,
 its actions agree on 358 queries; of 26 disagreements only eight change
 benchmark hit status, four in each direction. Relative to the rule
 there are 13 fixes and 14 harms. The fit IDs and H query IDs are disjoint,
 but the design of this text arm was chosen after H outcomes had already
 been inspected. A tie in this exposed-cohort outcome is not evidence
 of equivalence or a matched-cost performance guarantee. The parser is
 lossy despite lossless input serialization, and it is neither an
 LLM-baseline evaluation nor an adapted-I semantic-encoder control.

 \paragraph{Why input equality has several meanings.}
 We distinguish the underlying fact set, its textual or graph-oriented
 representation, model supervision, candidate pool, upstream ranking
 calls, and measured deployment resources. The original text12 ablation
 shares a candidate pool, rankings, semantic scores and label budget,
 but does \emph{not} have the link and edge facts. The post hoc parser
 has those facts, but is a deterministic program rather than a matched
 token/latency LLM. The H structured-text learner has the facts and
 original fit labels, but was designed after exposure and uses lossy
 lexicalization; there is no corresponding I adapted-text same-fact
 learner. The top-link rule has underlying facts but differs in features
 and support-first action logic. Conflating these controls would turn
 one valid ablation result into an invalid graph-exclusivity claim.

 \subsection{Measured local resources versus deployment cost}\label{app:resources}
 Two separate exploratory benchmarks reused \emph{exposed H}; neither read
 new query answers or assessed quality on a new cohort. The first used
 seven serial repeated predictions over 384 queries and 7,680 candidates,
 alternating graph and text arms with one prediction thread. Its graph
 input features had already been cached. The text arm separately
 serialized relation facts into JSON and parsed its 18 columns before
 model inference. Both models reproduced their original H actions on
 each repeat. Table~\ref{tab:resource} separates these stages; timings
 are descriptive local measurements, not confidence intervals or
 end-to-end deployment comparisons.

 \begin{table}[h]
 \centering
 \caption{Exposed-H local resource observations, seven runs per stage.
 Different protocols must not be compared across the panel boundary.}
 \label{tab:resource}\small
 \begin{tabular}{p{8.7cm}rr}
 \toprule
 Stage and measurement boundary & Median (s) & Range (s) \\
 \midrule
 Text JSON serialization and parse; no index & 0.550701 & 0.547480--0.756351 \\
 Graph model prediction; graph features cached & 0.023100 & 0.022765--0.030991 \\
 Text model prediction; text features parsed & 0.023335 & 0.023135--0.031144 \\
 \midrule
 Graph six-column construction, \emph{preloaded} facts & 0.126542 & 0.125997--0.128655 \\
 Text JSON and 18-column construction, \emph{same preloaded} facts & 0.113578 & 0.112287--0.114158 \\
 \bottomrule
 \end{tabular}
 \end{table}

 The second benchmark independently and symmetrically constructs graph
 and structured-text columns from the \emph{same preloaded} H case and
 link-option data, seven serial repeats with alternating arm order. Its
 graph producer exactly recovers the 46,080 frozen graph scalars on every
 repeat. The first benchmark includes Python \texttt{tracemalloc} while
 timing text serialization; it also records a 4,311,096-byte median
 Python allocation peak during that parse, \emph{not} process RSS.
 It is invalid to compare its 0.550701-second text parse directly with
 the other run's 0.126542-second graph construction as an efficiency
 ratio. The construction outputs also have different feature dimensions.

 Both benchmarks omit full-corpus index build/update, BM25/E5 retrieval,
 LLM reference generation and two ranking requests, entity-linking
 execution, Qwen scoring, model loading, process/VRAM peak and controlled
 cold/warm end-to-end latency. The recorded service \texttt{cost} fields
 have unverified units; dollars cannot be reconstructed reliably from
 them. Machine resources were not exclusively reserved to certify
 same-time-budget deployment. Neither arm is shown to be cheaper or
 noninferior under an actual matched end-to-end resource budget.

 \subsection{Exploratory cross-task and negative-control evidence}\label{app:mag}
 The exploratory STaRK-MAG author task retrieves papers, not Amazon
 products, and neither Amazon model is transferred to it. On a selected
 1,024-query MAG TRAIN development sample, author-writes reranking
 hits 376, author-enhanced BM25 363 and a mismatched-support control
 299. Yet exact matching against the same index's author text field
 reproduces every graph-rule action on its 122 linkable queries. This
 supports a counterexample to a \emph{representation-exclusive}
 interpretation of that particular graph path; it is not independent
 cross-dataset confirmation of Amazon's learned fusion.

 Other MAG TRAIN pilots supply deliberately negative diagnostic evidence.
 On 259 disjoint TRAIN queries, an external-author-ID graph profile hits
 168, exact-name text 185 and BM25 153; these graph and text profiles are
 \emph{not} strict equal-fact representations. On 434 collaborator
 queries a graph path and its named-text comparator differ on only four
 top-1 actions. On 512 institution queries, graph and same-fact text
 prompts differ on only nine actions. Those latter two comparisons failed
 their prespecified label-blind disagreement gates; their answers were
 not opened and no hit rates for them are claimed. The prompts do not
 have exactly equal lengths even when they share facts and model access.
 By Eq.~\eqref{eq:disagreement}, any Hit@1 difference on their fixed
 cohorts is at most $4/434$ or $9/512$ respectively, regardless of
 which policy would benefit. We did not evaluate MAG official VAL/TEST
 or human queries for these paths.

\clearpage
\section{Full frozen comparison families}\label{app:families}
Tables~\ref{tab:fullh} and~\ref{tab:fulli} show every prespecified
comparator, not only the favorable rows in the main paper. The first column
identifies a different policy on the \emph{same} cohort: one model's raw
ranking, a single-model LambdaRank (LTR), a linked-gain single-model
postprocessor, rank aggregation, or a graph-support rule. These policies
can differ in minimal deployment cost even when evaluated on the same
cached candidates and service outputs. In particular, rank-only policies
need not run the extra Qwen scorer or reference-extraction path if deployed
independently. Each table uses a \emph{single} numerator (the frozen graph
method) and a query denominator fixed before result inspection. A negative
net count is a result, not a missing experiment.
\begin{table}[h]
\centering
\caption{Complete H family ($N=384$; primary: 18D graph fusion, 278 hits).
Fix/harm is primary-only hit/comparator-only hit. Adjusted $p$ uses all ten
one-sided paired comparisons; rounding is for display only. LTR and linked
postprocessors use the indicated single ranking.}
\label{tab:fullh}\small
\begin{tabular}{lrrrl}
\toprule
Comparator & Hits & Fix/harm & Net & Holm $p$ \\
\midrule
Text fusion (12D) & 264 & 15/1 & +14 & 0.002594 \\
Shuffled fusion (18D) & 263 & 17/2 & +15 & 0.003279 \\
RRF & 261 & 26/9 & +17 & 0.023952 \\
Borda & 261 & 26/9 & +17 & 0.023952 \\
Single-model LTR 0 & 270 & 18/10 & +8 & 0.554800 \\
Single-model LTR 1 & 273 & 15/10 & +5 & 1 \\
Single-model linked 0 & 273 & 17/12 & +5 & 1 \\
Single-model linked 1 & 277 & 14/13 & +1 & 1 \\
Max-world support & 276 & 14/12 & +2 & 1 \\
Top-link support & 279 & 13/14 & $-1$ & 1 \\
\bottomrule
\end{tabular}
\end{table}

The H training pool consists of A/B/F, not H labels; development on C/D/E
was used to choose the method. The 384 H queries were held out \emph{at
the time of that freeze}, not across the complete research history. RRF
and Borda coincide in this cohort, making their two rows two controls in
the fixed family but not independent empirical replications. The apparent
positive difference over maximum-world support is only two queries. The
top-link control beats the main graph method by one hit; its nonsignificant
comparison must not be reversed into an equivalence claim. Complete
query-level action and hit records reside under the frozen H results
archive (see Appendix~\ref{app:trace}).

\begin{table}[h]
\centering
\caption{Complete I family ($N=768$; primary: adapted 20D graph fusion,
568 hits). The older 18D text/shuffle rows do \emph{not} share the adapted
encoder's additional supervision, unlike adapted text/shuffle.}
\label{tab:fulli}\small
\begin{tabular}{lrrrl}
\toprule
Comparator & Hits & Fix/harm & Net & Holm $p$ \\
\midrule
Old 18D graph fusion & 563 & 17/12 & +5 & 1 \\
Old text fusion & 542 & 33/7 & +26 & 0.000296 \\
Old shuffled fusion & 543 & 33/8 & +25 & 0.000729 \\
RRF & 550 & 37/19 & +18 & 0.111207 \\
Borda & 550 & 37/19 & +18 & 0.111207 \\
Single-model LTR 0 & 553 & 43/28 & +15 & 0.383695 \\
Single-model LTR 1 & 562 & 34/28 & +6 & 1 \\
Single-model linked 0 & 555 & 42/29 & +13 & 0.538695 \\
Single-model linked 1 & 565 & 33/30 & +3 & 1 \\
Max-world support & 570 & 24/26 & $-2$ & 1 \\
Top-link support & 573 & 20/25 & $-5$ & 1 \\
Matched adapted text (14D) & 549 & 24/5 & +19 & 0.003277 \\
Adapted shuffled graph (20D) & 552 & 21/5 & +16 & 0.013717 \\
Original adapted-encoder score & 561 & 17/10 & +7 & 0.743366 \\
Raw adapted model ranking & 431 & 167/30 & +137 & $<10^{-20}$ \\
\bottomrule
\end{tabular}
\end{table}

Five I comparisons are positive after the specified adjustment: matched
adapted text, adapted shuffle, old text, old shuffle, and the standalone
semantic score. The last three are not substitutes for an equal-supervision
graph comparison. The raw adapted model ranking is an included, much
weaker control, not evidence of strong-baseline superiority. Neither RRF
nor either graph-support rule is among the five positive adjusted
comparisons. The adapted model is not significantly better than the
old 18D graph model on I, so the additional encoder-training step itself
has no isolated confirmed gain. This family cannot be pooled with H's
ten tests or interpreted as 1,152 repeated trials of one policy.

\section{Derivations and finite-sample scope}\label{app:derivations}
\paragraph{Error accounting.}
For a policy required to select from $E_q\subseteq C_q$, the three terms in
Eq.~\eqref{eq:decomposition} are nonnegative because
$H_q(\pi)\le O_q(E_q)\le O_q(C_q)\le1$. Summing the three brackets
telescopes to $1-H_q(\pi)$. Neither $O_q(E_q)$ nor $O_q(C_q)$ can be
computed without the benchmark labels, and neither is an attainable
deployed policy. The full-pool oracle upper bound on hit counts is 354/384
for H and 704/768 for I. Thus the 18D graph fusion's 106 H misses include
30 without any labeled answer in the fixed pool and 76 with an answer
present but not selected. The adapted fusion's 200 I misses likewise
split into 64 unavailable and 136 within-pool selection failures. This
does not identify which failures arise specifically from the graph.

\paragraph{Paired outcomes.}
For each query define $d_q=H_q(g)-H_q(t)\in\{-1,0,1\}$. Its nonzero terms
are the $F$ fixes and $B$ harms, hence
$\sum_qd_q=F-B$. If two policies return the same candidate, both hit
indicators coincide for every possible $Y_q$, so $d_q=0$ whenever
$q\notin D(g,t)$. Thus
$|\sum_qd_q|\le\sum_{q\in D}|d_q|\le |D|$, proving
Eq.~\eqref{eq:disagreement}. This bound is about the \emph{observed
fixed queries} and policies, not a rate guarantee for future queries.
The paired exact test in Eq.~\eqref{eq:exact} conditions on the observed
number of discordant outcomes. It needs the interpretation that under the
null their two signs are exchangeable; Holm's family control applies to
those particular frozen tests, not to previous model search. A confidence
interval crossing zero does not certify noninferiority. To test that claim
one would need a justified margin and a separate prespecified procedure.

\paragraph{Matched-size, matched-rank-set expectation.}
Partition each fixed candidate pool into five consecutive bins of four
by the two-LLM RRF order. In bin $b$, let $n_b=4$ candidates include
$p_b$ benchmark positives, and let a support domain select $k_b$ items.
A control set uniformly selects $k_b$ without replacement in each bin.
Exactly $\binom{n_b-p_b}{k_b}$ of the $\binom{n_b}{k_b}$ possible choices
avoid positives. Since bin draws are independent by construction,
\begin{equation}\label{eq:hypergeom}
 \mathbb E[O_q(E_{\mathrm{control}})]
 =1-\prod_{b=1}^{5}
   \frac{\binom{n_b-p_b}{k_b}}{\binom{n_b}{k_b}}.
\end{equation}
The usual convention assigns zero to a bin's avoidance probability when
$k_b>n_b-p_b$. Equation~\eqref{eq:hypergeom} is an exact expectation for
this artificial matched-set distribution, not a model of how graph edges
are generated. A more finely matched degree or within-bin rank control
might change the descriptive comparison. On the \emph{post hoc} activated
subsets ($|E_q|<20$), recorded graph-domain oracles total 43/55 H and
79/94 I; the corresponding matched-set oracle expectations sum to 25.50
and 44.52. Fractional values are expectations, not fractional queries or
additional independently sampled observations. On those same subsets
top-link hits 37 and 67, while the matched-set first-position expectations
are 16.92 and 29.67. All these numbers were computed after H/I exposure;
no new confirmatory $p$ value follows.

\paragraph{A restricted improvement ceiling.}
Consider a hypothetical new policy that differs from top-link only on
the activated set $\mathcal A$ and still selects from $E_q$ there. For
each activated $q$, $H_q(\pi)\le O_q(E_q)$, while outside $\mathcal A$
the two policies have equal hits. Consequently its \emph{observed}
improvement is bounded by
\begin{equation}
 \sum_q[H_q(\pi)-H_q(\pi_{\mathrm{top}})]
 \le\sum_{q\in\mathcal A}
 [O_q(E_q)-H_q(\pi_{\mathrm{top}})].
\end{equation}
The right-hand side is six H queries and twelve I queries, each 1.5625
percentage points of its respective full cohort. These are label-oracle
ceilings conditional on the specified domain and activation restriction;
they do not bound policies that change choices outside $\mathcal A$ or
select outside $E_q$, and they do not promise an attainable gain. For
example, the studied full-pool fusion is \emph{not} constrained by the
restriction. No novel generalization theorem is asserted here.